% Generated draft: requires mathematical and source review.
\documentclass[11pt]{article}
\usepackage[margin=25mm]{geometry}
\usepackage{fontspec}
\setmainfont{Latin Modern Roman}
\usepackage{amsmath,amssymb,mathtools,unicode-math}
\setmathfont{Latin Modern Math}
\usepackage{enumitem,needspace,xurl,hyperref}
\setlength{\parindent}{0pt}
\setlength{\parskip}{5pt}
\setlength{\emergencystretch}{3em}
\newcommand{\N}{\mathbb{N}}
\newcommand{\Z}{\mathbb{Z}}
\newcommand{\Q}{\mathbb{Q}}
\newcommand{\R}{\mathbb{R}}
\newcommand{\C}{\mathbb{C}}
\newcommand{\F}{\mathbb{F}}
\newcommand{\A}{\mathbb{A}}
\newcommand{\PP}{\mathbb{P}}
\newcommand{\E}{\mathbb{E}}
\newcommand{\eps}{\varepsilon}
\newcommand{\dd}{\,\mathrm{d}}
\begin{document}
\section*{UnsolvedMath ID: 6700061}
\textbf{Source problem code:} catalogue entry
\subsection{Definitions and mathematical statement}
Let $\Delta\subset\R^3$ be a regular Euclidean tetrahedron, with corresponding faces and edges fixed. In the interior-dihedral convention its common angle is $\theta_0=\arccos(1/3)$. The source asks whether $\Delta$ is mean-convexly extremal, by asking about a competing simplex $P$ mapped facewise to $\Delta$ without decreasing distances.

Under a face-metric reading, the distance requirement would be
\[
d_{\Delta_i}(f_i(x),f_i(y))\geq d_{P_i}(x,y)
\qquad(x,y\in P_i)
\]
for each corresponding face map $f_i:P_i\to\Delta_i$. This retains the stated nondecreasing direction; it is not replaced by distance contraction. The competitor is to have nonnegative face mean curvatures and dihedral angles no larger than the corresponding angles of $\Delta$, with at least one strictly smaller.

The area-extremality-style question is therefore whether such a strict-angle competitor can exist. Affirmative mean-convex extremality would rule it out; a yes answer to the existence question would refute that extremality. The source's word ``equivalently'' joins opposite yes/no formulations, not identical affirmative assertions.

The competitor category is not defined: an ordinary straight Euclidean simplex, a curvilinear Euclidean domain, and a Riemannian simplex with corners give different problems. The source supplies no interior scalar-curvature condition, and none is invented. The face distance metric, face-map regularity/bijectivity, mean-curvature normal/sign, and interior versus exterior dihedral convention must be specified. For curved edges, the meaning of one angle being strictly smaller also needs a pointwise or edgewise convention. The source does not demand that every non-strict competitor be isometric to $\Delta$; that stronger rigidity claim is not added.
\subsection{Short English statement}
Can a facewise distance-nondecreasing comparison to the regular tetrahedron have mean-convex faces and no larger dihedral angles, with one angle smaller? Mean-convex extremality would forbid such a competitor, subject to the missing geometric conventions.
\subsection{Sources}
\begin{itemize}
\item Local UnsolvedMath registry, permanent ID 6700061, exact source code catalogue entry; proposed by Misha Gromov, 2017. Registry status open, classification OPEN-TRIAGE. Source printed tag [?65], occurrence 64.
\item Misha Gromov, \emph{101 Questions, Problems and Conjectures around Scalar Curvature} (2017), printed tag [?65], PDF page 61, \url{https://www.ihes.fr/~gromov/wp-content/uploads/2018/08/101-problemsOct1-2017.pdf}. Supplied normalized pdf-text extraction accessed July 29, 2026; incomplete author-hosted PDF. No PDF was accessed here.
\item The supplied summary calls this a concrete mean-curvature/dihedral question and reports no resolution located. It supplies no precise later reference or missing definition of mean-convex extremality.
\item Provenance flags current status and statement-reuse terms NEEDS\_REVIEW; the source marks the passage as a numbered open item.
\end{itemize}
\end{document}
